\documentclass{article}
\usepackage[utf8]{inputenc}
\usepackage{a4wide}

\begin{document}

\section*{Rider น้อย}

ในเวลาว่างณัฏฐกิตติ์จะรับงานเสริมเป็น rider

ที่ต้องขี่มอเตอร์ไซค์ไปส่งอาหารจากร้านอาหารเขารับออเดอร์ซึ่งตั้งอยู่ที่จุด 0 บนเส้นจำนวน

ไปยังบ้านของลูกค้าที่จุด d บนเส้นจำนวน มอเตอร์ไซค์ของเขามีถังน้ำมันที่บรรจุได้ n ลิตร

และถังน้ำมันจะเต็มเมื่อเริ่มต้นการเดินทาง โดยในการเดินทาง 1 หน่วยบนเส้นจำนวนจะใช้น้ำมัน

1 ลิตร

ในขณะขี่มอเตอร์ไซค์เดินทางณัฏฐกิตติ์จะต้องหยุดเติมน้ำมันที่สถานีบริการน้ำมันที่ตั้งอยู่ตามจุดต่าง ๆ

บนเส้นทาง โดยสถานีแต่ละแห่งจะมีราคาน้ำมันที่แตกต่างกัน และน้ำมันที่ใช้จะหมดไป 1

ลิตรทุกระยะทาง 1 หน่วย



\ul{งานของคุณ}: ช่วยหาว่าในการเดินทางณัฏฐกิตติ์จะเสียค่าน้ำมันน้อยที่สุดเท่าใด

เพื่อให้เขาสามารถเดินทางไปส่งออเดอร์ได้สำเร็จ



\hfill\break



{\def\LTcaptype{none} % do not increment counter

\begin{longtable}[]{@{}

  >{\raggedright\arraybackslash}p{(\linewidth - 4\tabcolsep) * \real{0.3333}}

  >{\raggedright\arraybackslash}p{(\linewidth - 4\tabcolsep) * \real{0.3333}}

  >{\raggedright\arraybackslash}p{(\linewidth - 4\tabcolsep) * \real{0.3333}}@{}}

\toprule\noalign{}

\begin{minipage}[b]{\linewidth}\centering

ข้อมูลเข้า

\end{minipage} & \begin{minipage}[b]{\linewidth}\centering

ข้อมูลออก

\end{minipage} & \begin{minipage}[b]{\linewidth}\centering

คำอธิบาย

\end{minipage} \\

\midrule\noalign{}

\endhead

\bottomrule\noalign{}

\endlastfoot

\begin{minipage}[t]{\linewidth}\raggedright

10 4 4\\

3 5\\

5 8\\

6 3\\

8 4\\

\strut

\end{minipage} & 22 & \begin{minipage}[t]{\linewidth}\raggedright

ระยะทางที่ต้องเดินทางคือ 10 หน่วย\\

มอเตอร์ไซค์สามารถเดินทางได้ไกลที่สุด 4 หน่วยโดยไม่ต้องเติมน้ำมัน\\

ณัฏฐกิตติ์จะเดินทางไปสถานีบริการที่จุด 3 (เติม 2 ลิตร, ราคา 5\\

หน่วยต่อลิตร) แล้วเดินทางต่อไปที่สถานีบริการที่จุด 6 (เติม 4 ลิตร, ราคา 3\\

หน่วยต่อลิตร) และสุดท้ายจะไปถึงบ้านลูกค้า ค่าทั้งหมดที่จ่ายคือ 2 * 5 + 4 * 3 = 22 หน่วย\\

\strut

\end{minipage} \\

\begin{minipage}[t]{\linewidth}\raggedright

16 5 2\\

8 2\\

5 1\\

\strut

\end{minipage} & -1 & \begin{minipage}[t]{\linewidth}\raggedright

ระยะทางที่ต้องเดินทางคือ 16 หน่วย\\

มอเตอร์ไซค์สามารถเดินทางได้ไกลที่สุด 5 หน่วย\\

ซึ่งหมายความว่าณัฏฐกิตติ์สามารถขับรถไปได้ไกลสุดแค่ 5\\

หน่วยก่อนน้ำมันจะหมด จึงต้องเติมน้ำมันที่สถานีบริการที่จุด 5\\

และไปเติมต่อที่จุด 8 แต่ระยะทางถึงบ้านลูกค้าคือ 8 หน่วย\\

เกินกว่าที่ณัฏฐกิตติ์จะขี่มอเตอร์ไซค์ไปได้\\

ดังนั้นเขาจึงไม่สามารถไปถึงบ้านลูกค้าได้ ผลลัพธ์คือ -1\\

\strut

\end{minipage} \\

\end{longtable}

}



\subsection*{Input}
บรรทัดที่ 1: จำนวนเต็ม 3 ค่า แยกด้วยช่องว่าง ได้แก่ d (ระยะทางจากร้านไปยังบ้านลูกค้า),

n (ขนาดของถังน้ำมัน) , m (จำนวนสถานีบริการน้ำมัน) โดยที่ 1 ≤ n ≤ d ≤ 10\^{}9,

1 ≤ m ≤ 200,000



บรรทัดที่ 2 เป็นต้นไป จำนวน m บรรทัด: จำนวนเต็ม 2 ค่า แยกด้วยช่องว่าง ต่อ 1 บรรทัด

ได้แก่ xi (ตำแหน่งของสถานีบริการ) และ pi (ราคาน้ำมันต่อลิตร) โดยที่

1 ≤ x{[}i{]} ≤ d - 1, 1 ≤ p{[}i{]} ≤ 10\^{}6



\subsection*{Output}
บรรทัดที่ 1: ค่าน้ำมันน้อยที่สุดที่สามารถเดินทางไปส่งออเดอร์ได้สำเร็จ

ถ้าหากไม่สามารถเดินทางไปส่งออเดอร์สำเร็จ ให้แสดงผลเป็น -1\\



\end{document}
